\subsection{拓扑空间的附着}

设 X、B 为拓扑空间，设 A 为 B 的子空间，并设 \(f: A \rightarrow X\) 为连续映射。用 Y 表示空间 \(Y_{1} = X\) 与 \(Y_{2} = B\) 的拓扑和，并借助典范单射把 \(Y_{1}\) 与 \(Y_{2}\) 等同于 Y 的子空间。设 R 为 Y 上最细的等价关系，使得 \(Y_{2}\) 的元素 a 与 \(f(a)\)（即 \(Y_{1}\) 的元素）等价（对每个 \(a \in A\)）。关系 R 在 \(Y_{1}\) 上就是相等关系。设 \(x \in Y_{1}\) 与 \(b \in Y_{2}\)；我们有 x R b 当且仅当 \(b \in A\) 且 \(f(b) = x\)。设 \(b, b' \in Y_{2}\)；要使 b R \(b'\) 成立，必须而且只需 \(b = b'\)，或者 \(b, b' \in A\) 且 \(f(b) = f(b')\)。

定义 9. --- 商拓扑空间 Y/R 称为沿映射 f 把空间 B 附着到空间 X 上所得的空间，记作 \(X \cup_{f} B\)。

用 \(\alpha_{f}\colon\mathrm{X}\to\mathrm{X}\cup_{f}\mathrm{B}\) 与 \(\beta_{f}\colon\mathrm{B}\to\mathrm{X}\cup_{f}\mathrm{B}\) 表示 Y 到 Y/R 上的典范满射在 \(\mathrm{Y}_{1}\) 与 \(\mathrm{Y}_{2}\) 上的限制。我们有 \(\alpha_{f}\circ f=\beta_{f}\mid\mathrm{A}\)。

注记。--- 1) 映射 \(\alpha_{f}\) 是单射。对 X 的每个子集 U，有 \(\bar{\alpha}_{f}^{-1}(\alpha_{f}(U)) = U\) 与 \(\bar{\beta}_{f}^{-1}(\alpha_{f}(U)) = f^{-1}(U)\)。若 A 是 B 的开（相应地，闭）子集，这些关系式蕴涵 \(\alpha_{f}\) 是从 X 到 \(X \cup_{f} B\) 中的开（相应地，闭）映射。

设 U 是 X 的开子集，V 是 B 的开集，且满足 \(f^{-1}(U) = V \cap A\)。\(Y_{1}\) 的子集 U 与 \(Y_{2}\) 的子集 V 的并是 Y 的开饱和子集；因此它在 \(X \cup_{f} B\) 中的像是开集，它在 \(\alpha_{f}(X)\) 上的迹是 \(\alpha_{f}(U)\)。于是，映射 \(\alpha_{f}\) 定义了 X 到它在 \(X \cup_{f} B\) 中的像上的一个同胚。

\begin{enumerate}
\def\labelenumi{\arabic{enumi})}
\setcounter{enumi}{1}
\tightlist
\item
  设 V 是 B 的一个子集。我们有 \(\bar{\alpha}_{f}^{-1}(\beta_{f}(V)) = f(V \cap A)\) 与 \(\bar{\beta}_{f}^{-1}(\beta_{f}(V)) = V \cup \bar{f}(f(V \cap A))\)。若 A 在 B 中闭且映射 f 是闭映射，则映射 \(\beta_{f}\) 是闭映射。类似地，若 A 在 B 中开且映射 f 是开映射，则映射 \(\beta_{f}\) 是开映射。
\end{enumerate}

在这两种情形下，映射 \(\beta_{f}\) 都诱导出 \(B\setminus A\) 到其像上的一个同胚。

\begin{enumerate}
\def\labelenumi{\arabic{enumi})}
\setcounter{enumi}{2}
\tightlist
\item
  假设 A 是闭的且映射 \(f\) 是逆紧的。我们刚才已看到映射 \(\beta_f\) 是闭映射。对每个 \(x \in X\)，有 \(\overline{\beta_f}(\alpha_f(x)) = \overline{f}(x)\)。因此它是 A 的一个拟紧子集（TG, I, p. 75, 定理 1），从而也是 B 的拟紧子集。对每个 \(b \in B\)，\(\overline{\beta_f}(\beta_f(b))\) 等于 \(\{b\}\)（若 \(b \in B - A\)），或等于 \(\overline{f}(f(b))\)（若 \(b \in A\)）；两种情形下它都是 B 的拟紧子集。因此映射 \(\beta_{f}\) 的纤维都是拟紧的；从而（同前所引）映射 \(\beta_{f}\) 是逆紧的。
\end{enumerate}

映射 \(\alpha_{f}\) 也是逆紧的（TG, I, p. 72, 命题 2）。由此推出，从 Y 到 X \(\cup_{f}\) B 上的典范映射是逆紧的。

例。--- 1) 设 X 与 Y 为拓扑空间，\(f\colon X\to Y\) 为连续映射。柱 \(\operatorname{Cyl}(f)\) 无非就是把空间 \(X\times I\) 沿映射 \(f\circ \mathrm{pr}_1\)（从 \(X\times \{1\}\) 到 Y 中）附着到空间 Y 上所得的空间。

\begin{enumerate}
\def\labelenumi{\arabic{enumi})}
\setcounter{enumi}{1}
\tightlist
\item
  设 X 为拓扑空间，B 为拓扑空间，A 为 B 的子空间，并设 \(f: A \to X\) 为连续映射，它诱导出 A 到其像上的一个同胚。
\end{enumerate}

令 \(\mathrm{A}' = f(\mathrm{A})\)；设 \(f'\) 是从 \(\mathrm{A}'\) 到 B 中的映射，它把每个 \(x \in \mathrm{A}'\) 对应到 \(x\) 在 \(f\) 下的唯一原像。映射 \(f' \circ f\) 是连续的；由于 \(f\) 是严格映射，映射 \(f'\) 是连续的。存在唯一的连续映射 \(u\)，从空间 \(\mathrm{X} \cup_f \mathrm{B}\) 到空间 \(\mathrm{B} \cup_{f'} \mathrm{X}\)，它把 \(\alpha_f(x)\) 映为 \(\beta_{f'}(x)\)（\(x \in \mathrm{X}\)），把 \(\beta_f(b)\) 映为 \(\alpha_{f'}(b)\)（\(b \in \mathrm{B}\)）；它是一个同胚，把子空间 \(\alpha_f(\mathrm{X})\) 映到子空间 \(\beta_{f'}(\mathrm{X})\) 上。

命题 8. --- 设 Z 为拓扑空间。

\begin{enumerate}
\def\labelenumi{\alph{enumi})}
\item
  设 \(u \colon X \to Z\) 与 \(v \colon B \to Z\) 为连续映射，满足 \(v(a) = u(f(a))\)（对所有 \(a \in A\)）。于是存在唯一的连续映射 \(w \colon X \cup_f B \to Z\)，使得 \(w \circ \alpha_f = u\) 且 \(w \circ \beta_f = v\)。
\item
  设 \(\sigma\colon\mathbf{X}\times\mathbf{I}\to\mathbf{Z}\) 与 \(\tau\colon\mathbf{B}\times\mathbf{I}\to\mathbf{Z}\) 为同伦。假设 \(\tau\) 在 A 上是固定的，且 \(\sigma(f(a),t)=\tau(a,t)\)（对所有 \(a\in\mathbf{A}\) 及每个 \(t\in\mathbf{I}\)）。于是存在唯一的同伦 \(\eta\colon(\mathbf{X}\cup_{f}\mathbf{B})\times\mathbf{I}\to\mathbf{Z}\)，使得 \(\eta(\alpha_{f}(x),t)=\sigma(x,t)\) 与 \(\eta(\beta_{f}(b),t)=\tau(b,t)\)（对 \(x\in\mathbf{X}\)、\(b\in\mathbf{B}\) 与 \(t\in\mathbf{I}\)）。这个同伦在 \(\beta_{f}(\mathbf{A})\) 上是固定的。
\end{enumerate}

本命题由商空间的定义及命题 10（I, p. 19）立即得出。

推论 1. --- 设 \(B'\) 为拓扑空间，\(A'\) 为 \(B'\) 的子空间，\(f': A' \to X\) 为连续映射。设 \(v: B \to B'\) 为连续映射，满足 \(v(A) \subset A'\) 且 \(f = f' \circ (v \mid A)\)。设 \(w: X \cup_f B \to X \cup_{f'} B'\) 为唯一的连续映射，满足 \(w \circ \alpha_f = \alpha_{f'}\) 且 \(w \circ \beta_f = \beta_{f'} \circ v\)。若 \(v\) 定义了偶 \((B, A)\) 到偶 \((B', A')\) 上的同伦等价，则 \(w\) 定义了偶 \((X \cup_f B, \alpha_f(X))\) 到偶 \((X \cup_{f'} B', \alpha_{f'}(X))\) 上的同伦等价。

设 \(v' \colon \mathrm{B}' \to \mathrm{B}\) 为连续映射，\(\tau \colon \mathrm{B} \times \mathbf{I} \to \mathrm{B}\) 为在 A 上固定、连接 \(\mathrm{Id}_{\mathrm{B}}\) 与 \(v' \circ v\) 的同伦，\(\tau' \colon \mathrm{B}' \times \mathbf{I} \to \mathrm{B}'\) 为在 \(\mathrm{A}'\) 上固定、连接 \(\mathrm{Id}_{\mathrm{B}'}\) 与 \(v \circ v'\) 的同伦。对 \(a' \in \mathrm{A}'\) 与 \(a = v'(a')\)，我们有 \(a' = v(a)\) 且 \(\beta_f(v'(a')) = \beta_f(a) = \alpha_f(f(a)) = \alpha_f(f'(a'))\)。由命题 8(a)，存在唯一的映射 \(w' \colon X \cup_{f'} B' \to X \cup_f B\)，满足 \(w' \circ \alpha_{f'} = \alpha_{f'}\) 且 \(w' \circ \beta_{f'} = \beta_f \circ v'\)。由命题 8(b)，存在唯一的同伦 \(\eta \colon (\mathrm{X} \cup_f \mathrm{B}) \times \mathbf{I} \to X \cup_{f'} B'\)，满足 \(\eta(\alpha_f(x), t) = \alpha_{f'}(x)\) 与 \(\eta(\beta_f(b), t) = \beta_{f'}(\tau(b, t))\)（对 \(x \in X\)、\(b \in B\) 与 \(t \in I\)）。同样存在唯一的同伦 \(\eta' \colon (X \cup_{f'} B') \times I \to X \cup_f B\)，满足 \(\eta'(\alpha_{f'}(x), t) = \alpha_f(x)\) 与 \(\eta'(\beta_{f'}(b), t) = \beta_f(\tau'(b, t))\)（对 \(x \in X\)、\(b \in B'\) 与 \(t \in I\)）。于是，从 \((X \cup_f B) \times I\) 到 \(X \cup_f B\) 中的映射 \((x, t) \mapsto \eta'(\eta(x, t), t)\) 是在 \(\alpha_f(X)\) 上固定、连接 \(X \cup_f B\) 的恒同映射与映射 \(w' \circ w\) 的同伦。类似地，从 \((X \cup_{f'} B') \times I\) 到 \(X \cup_{f'} B'\) 中的映射 \((x, t) \mapsto \eta(\eta'(x, t), t)\) 是在 \(\alpha_{f'}(X)\) 上固定、连接 \(X \cup_{f'} B'\) 的恒同映射与映射 \(w \circ w'\) 的同伦。推论由此得出。

例 3. --- 用 i 表示 A 到 B 中的典范单射，取映射柱 i 作为空间 B'；用 \(\alpha_{i}\colon\mathrm{A}\times\mathbf{I}\to\mathrm{Cyl}(i)\) 与 \(\beta_{i}\colon\mathrm{B}\to\mathrm{Cyl}(i)\) 表示典范映射，用 \(r_{i}\colon\mathrm{Cyl}(i)\to\mathrm{B}\) 表示柱 \(\mathrm{Cyl}(i)\) 到其柱底上的典范收缩映射。设 \(A_{0}\) 为 \(\alpha_{i}(\mathrm{A}\times\{0\})\)，即 \(\mathrm{Cyl}(i)\) 的子空间，并用 \(f_{0}\colon\mathrm{A}_{0}\to\mathrm{X}\) 表示映射 \(f\circ r_{i}\)。对 \(a\in A\)，我们有

\[
\beta_ {f} \circ r _ {i} (\alpha_ {i} (a, 0)) = \beta_ {f} (a) = \alpha_ {f} (f _ {0} (\alpha_ {i} (a, 0))) = \alpha_ {f} (f (a)).
\]

设 \(\eta\colon X\cup_{f_0}\mathrm{Cyl}(i)\to X\cup_fB\) 为唯一的连续映射，满足 \(\eta\circ\alpha_{f_0}=\alpha_f\) 且 \(\eta\circ\beta_{f_0}=\beta_f\circ r_i\)。假设偶 (B, A) 具有同伦扩张性质。于是映射 \(r_i\) 定义了偶 \((\mathrm{Cyl}(i),\mathrm{A}_0)\) 到偶 (B,A) 上的同伦等价（III, p. 243, 定理 1）。由推论 1，映射 \(\eta\) 便定义了偶 \((X\cup_{f_0}\mathrm{Cyl}(i),\alpha_{f_0}(X))\) 到偶 \((X\cup_fB,\alpha_f(X))\) 上的同伦等价。特别地，\(\eta\) 是一个同伦等价。

推论 2. --- 设 \(X'\) 为拓扑空间，\(u: X \to X'\) 为连续映射，并令 \(f' = u \circ f\)。设 \(w: X \cup_f B \to X' \cup_{f'} B\) 为唯一的连续映射，满足 \(w \circ \alpha_f = \alpha_{f'} \circ u\) 且 \(w \circ \beta_f = \beta_{f'}\)。若 \(u\) 定义了偶 \((X, f(A))\) 到偶 \((X', f'(A))\) 上的同伦等价，则 \(w\) 定义了偶 \((\mathrm{X}\cup_f\mathrm{B},\beta_f(\mathrm{B}))\) 到偶 \((\mathrm{X}'\cup_{f'}\mathrm{B},\beta_{f'}(\mathrm{B}))\) 上的同伦等价。

证明与推论 1 的证明类似。

命题 9. --- 设 X 与 B 为拓扑空间，A 为 B 的子空间，f: A → X 为连续映射。

\begin{enumerate}
\def\labelenumi{\alph{enumi})}
\item
  若偶 (B, A) 具有同伦扩张性质，则偶 \((\mathrm{X} \cup_{f} \mathrm{B}, \alpha_{f}(\mathrm{X}))\) 亦然。
\item
  假设映射 f 是单射且严格的。若偶 \((\mathrm{X}, f(\mathrm{A}))\) 具有同伦扩张性质，则偶 \((\mathrm{X} \cup_{f} \mathrm{B}, \beta_{f}(\mathrm{B}))\) 亦然。
\item
  假设偶 (B, A) 具有同伦扩张性质。设 Z 为拓扑空间，\(u: X \cup_{f} B \to Z\) 为连续映射，\(\sigma: \alpha_{f}(X) \times I \to Z\) 为一同伦，其末项映射是 u 在子空间 \(\alpha_{f}(X)\) 上的限制。
\end{enumerate}

令 \(v_{1} = u \circ \alpha_{f}\)，并用 \(\tau_{1} \colon \mathbf{X} \times \mathbf{I} \to \mathbf{Z}\) 表示由 \((x,t) \mapsto \sigma(\alpha_{f}(x),t)\) 给出的映射。令 \(v_{2} = u \circ \beta_{f}\)。由于 \(\beta_{f} \mid \mathrm{A} = \alpha_{f} \circ f\)，从 \(\mathrm{A} \times \mathbf{I}\) 到 \(\mathrm{Z}\) 中、把 \((a,t)\) 映到 \(\sigma(\alpha_{f}(f(a)),t)\)（\(a \in \mathrm{A}\)，\(t \in \mathbf{I}\)）的映射是一个同伦，其末项等于映射 \(v_{2} \mid \mathrm{A}\)。由于偶 \((\mathrm{B},\mathrm{A})\) 具有同伦扩张性质，存在同伦 \(\tau_{2} \colon \mathbf{B} \times \mathbf{I} \to \mathbf{Z}\)，其末项为映射 \(v_{2}\)，且满足 \(\tau_{2}(a,t) = \sigma(\alpha_{f}(f(a)),t)\)（\((a,t) \in \mathrm{A} \times \mathbf{I}\)）。

对 \(a \in \mathbf{A}\) 与 \(t \in \mathbf{I}\)，我们有 \(\tau_2(a, t) = \tau_1(f(a), t)\)。由命题 8，存在唯一的连续映射 \(\sigma: (\mathrm{X} \cup_f \mathrm{B}) \times \mathbf{I} \to \mathbf{Z}\)，满足 \(\sigma(\alpha_f(x), t) = \sigma_1(x, t)\) 与 \(\sigma(\beta_f(b), t) = \sigma_2(b, t)\)（对 \(x \in \mathrm{X}\)、\(b \in \mathrm{B}\) 与 \(t \in \mathbf{I}\)）。这是一个以 \(u\) 为末项、开拓 \(\tau\) 的同伦，由此得断言 a)。

\begin{enumerate}
\def\labelenumi{\alph{enumi})}
\setcounter{enumi}{1}
\tightlist
\item
  注意到例 2（III, p. 249），断言 b) 由断言 a) 得出。
\end{enumerate}

